% RECOVERY32 reconstruction from retained complete proof/context and restored31 inputs.
% Not asserted byte-identical to the lost pre-outage SHA9144ff4f... source.
\subsection*{The retained odd-boundary Schreier kernel and its actual relative pro-two topology}
\statement{A.102}{Proved}{Every marked odd boundary cover retains its full free pro-two Schreier kernel and the original canonical involution}
\label{stmt:A102}
Use the actual sphere quotient and original full-GT involution of
A.96--A.98:
\[
 G=N\rtimes B,\qquad N=F(P,A,C),\quad B=F(d,f),\qquad
 P=def,\quad A=af^{-1},\quad C=cd^{-1}.
\]
For an odd integer $n>1$ let
\[
 N_n=\ker\bigl(N\to C_n,\ P\mapsto1,\ A,C\mapsto0\bigr),
 \quad p=P^n,\quad A_i=P^iAP^{-i},\quad C_i=P^iCP^{-i}.
\]
Then $N_n$ has the literal free basis
\[
 p,\ A_0,\ldots,A_{n-1},\ C_0,\ldots,C_{n-1}.
\]
The four signed base actions are the full indexed formulas below, with
their nonperiodic $p$ wraps retained. There is an actual finite head
$G\to C_n\times D_{2n}$ with kernel $K_n=N_n\rtimes B_n$ for which
\[
 K_n^{(2)}=N_n^{(2)}\rtimes B_n^{(2)}.
\]
Consequently the actual relative quotient
\[
 G_{(n,\mathrm{rel})}
 =\widehat G/O^2(\widehat K_n)
\]
contains the embedded free pro-two group $N_n^{(2)}$ on the displayed
basis. Both the original involution $\alpha$ and its precisely
specified normalized reflection act on this quotient, with the full
correction $h_0$ retained inside that kernel. No pair-twist square or
class-two relation is imposed in this proposition.
\begin{proof}
\emph{The marked Schreier basis and every signed row.}
The transversal $1,P,\ldots,P^{n-1}$ gives the stated basis by ordinary
Schreier rewriting. For every integer $i$ use the exact convention
\begin{equation}\label{eq:A102-wrap}
 A_{i+n}=pA_ip^{-1},\qquad C_{i+n}=pC_ip^{-1}.
\end{equation}
In particular an index $-1$ is a conjugate by $p^{-1}$, and an index
$n$ is a conjugate by $p$; the indices are not simply reduced modulo
$n$ in the group.

Let $D=\operatorname{Ad}(d)$ and $F=\operatorname{Ad}(f)$, using
positive conjugation. The entire four signed actions are
\begin{equation}\label{eq:A102-signed}
\begin{aligned}
 D(p)&=C_0^{-1}pC_0,&D(C_i)&=C_0^{-1}C_iC_0,\\
 D(A_i)&=C_0^{-1}C_{i-1}A_iC_i^{-1}C_0,&&\\
 D^{-1}(p)&=C_0pC_0^{-1},&
 D^{-1}(C_i)&=C_0C_iC_0^{-1},\\
 D^{-1}(A_i)&=C_0C_{i-1}^{-1}A_iC_iC_0^{-1},&&\\
 F(p)&=A_0^{-1}pA_0,&F(A_i)&=A_0^{-1}A_iA_0,\\
 F(C_i)&=A_0^{-1}A_iC_iA_{i+1}^{-1}A_0,&&\\
 F^{-1}(p)&=A_0pA_0^{-1},&F^{-1}(A_i)&=A_0A_iA_0^{-1},\\
 F^{-1}(C_i)&=A_0A_i^{-1}C_iA_{i+1}A_0^{-1}.&&
\end{aligned}
\end{equation}
They follow from the actual A.96 rows
\[
 D(P)=C^{-1}PC,\quad D(A)=C^{-1}P^{-1}CPA,\quad D(C)=C,
\]
\[
 F(P)=A^{-1}PA,\quad F(A)=A,\quad F(C)=CPA^{-1}P^{-1}A.
\]
For the inverse $D$, the free basis $(P,C,PA)$ shows that $P$ is
conjugated by $C$ while $C,PA$ are fixed. For the inverse $F$, use
$(P,A,CP)$ and conjugate $P$ by $A$ while fixing $A,CP$.
Conjugating the resulting $A,C$ words by the images of $P^i$ gives
every row of \eqref{eq:A102-signed}. These arguments apply to negative
as well as positive $i$, so they retain every boundary wrap in
\eqref{eq:A102-wrap}. The $P$ exponent is preserved by all four maps;
in particular $N_n$ is invariant.

\emph{The complete mod-two head, not a selected eigenspace.}
Put $R_n=\mathbf F_2[T]/(T^n-1)$. The full elementary abelian quotient
of $N_n$ is
\[
 N_n/\Phi_2(N_n)
   =\mathbf F_2[p]\oplus R_n[A]\oplus R_n[C].
\]
The coordinate $[p]$ is fixed, and on the whole two-copy module the
actual action is
\begin{equation}\label{eq:A102-linear}
 D(a,c)=(a,c+(1+T^{-1})a),\qquad
 F(a,c)=(a+(1+T)c,c).
\end{equation}
These formulas follow by abelianizing every indexed row; the wraps
become periodic only at this abelian stage. Both operators square to
the identity. Their product has determinant one and characteristic
polynomial $(t+T)(t+T^{-1})$. Since $n$ is odd, $T^n-1$ is square-free.
After extending the finite field to split it, the $T=1$ component has
both operators equal to the identity, and every component $T=\zeta\ne1$
has distinct product eigenvalues $\zeta,\zeta^{-1}$. Thus the product
has order dividing $n$ on the ENTIRE module. A primitive $n$th root
gives exact order $n$. The reflection is nontrivial and has order two,
so the image is $D_{2n}$, not an assumed quotient detected only on one
character.

Let $B_n$ be the kernel of this actual map $B\to D_{2n}$. The base
acts trivially on $N/N_n=C_n$, so the resulting head of $G$ is exactly
$C_n\times D_{2n}$ and its kernel is
$K_n=N_n\rtimes B_n$. In particular $B_n$ acts trivially on the WHOLE
Frattini quotient of $N_n$, including $[p]$.

\emph{The full actual relative topology.}
A.96 supplies the full induced profinite topology of $N$ in $G$;
the finite-index subgroup $N_n$ has its full topology as well.
The same characteristic-quotient argument gives the actual completion
of the open subgroup
$\widehat K_n=\widehat N_n\rtimes\widehat B_n$.

Take any characteristic finite two-group quotient $J$ of the finitely
generated free group $N_n$. The image of $B_n$ in $\operatorname{Aut}(J)$
is trivial on $J/\Phi(J)$, by the complete module calculation above.
The kernel of
$\operatorname{Aut}(J)\to\operatorname{Aut}(J/\Phi(J))$ is a two-group.
Here is the needed elementary argument. In a finite $p$-group,
every proper subgroup is strictly contained in its normalizer, by
induction using the center. Hence maximal subgroups are normal of
index $p$, and $\Phi(J)=J^p[J,J]$. Every tuple lifting a basis of
$J/\Phi(J)$ generates $J$, since otherwise it lies in a maximal
subgroup. The automorphism kernel acts freely on the set of all such
lifts: fixing a generating tuple means fixing the group. The number
of lifts is $|\Phi(J)|^r$, a $p$-power, so the kernel order divides
a $p$-power. Apply this with $p=2$.

Therefore $B_n$ has a finite two-group action image on $J$.
Intersecting this action kernel with the kernel of any prescribed
finite two-group target of $B_n$ gives a finite two-group
semidirect target of $K_n$ retaining $J$ and that base target.
Characteristic finite two-group kernels in $N_n$ are cofinal:
there are only finitely many kernels of a bounded index, and their
automorphic intersection is still a finite two-group kernel.
These compatible targets yield the two inverse completion maps and
prove
\[
 K_n^{(2)}=N_n^{(2)}\rtimes B_n^{(2)}.
\]
Equivalently $O^2(\widehat K_n)\cap\widehat N_n=O^2(\widehat N_n)$.
Since $\widehat K_n$ is normal of finite index in $\widehat G$, its
characteristic pro-two residual is normal in $\widehat G$.
The relative quotient therefore retains precisely the asserted
embedded $N_n^{(2)}$, with primitive $p$ and the literal wraps.
The ENTIRE finite-head kernel is $K_n^{(2)}$, not just the retained
fiber $N_n^{(2)}$.

\emph{The original involution and its whole correction.}
The actual labelled peripheral proof in A.97 gives
$\alpha(\widehat N)=\widehat N$, action $-I$ on its abelianization,
and a normalized reflection
$\beta=\operatorname{Inn}(h_0)\alpha$ with
\[
 \beta|_{\widehat B}=\alpha_0,\qquad
 \alpha(b)=h_0^{-1}\alpha_0(b)h_0 .
\]
A.98 gives the exact integral coordinate
$(h_0)_{\mathrm{ab}}=-A$. Thus its $P$ coordinate is zero, so
$h_0\in\widehat N_n$ for EVERY odd $n$. The full correction is not
erased or replaced by an arbitrary coboundary.

The negative action on $N_{\mathrm{ab}}$ preserves the kernel of its
$P$ character modulo $n$, so both $\alpha$ and $\beta$ preserve
$\widehat N_n$. On the actual base, $\alpha_0$ sends each of
$d,f,df$ to a conjugate of its inverse. Hence it preserves the closed
normal kernel generated by $d^2,f^2,(df)^n$, which is exactly the
displayed dihedral kernel $\widehat B_n$. The normalized $\beta$
therefore preserves $\widehat K_n$, and the original $\alpha$ does
as well, since their difference is conjugation by the actual
$h_0\in\widehat N_n$. They preserve $O^2(\widehat K_n)$ and descend
to the stated relative quotient with the same full correction.
\end{proof}
The reconstruction uses the restored complete proofs
\path{work_A/actual_sphere_kernel_relative_two31.md},
\path{work_A/actual_sphere_canonical_reflection31.md}, and
\path{work_A/actual_sphere_cabling_correction31.md}, together with
the retained full cycle32 indexed proof. It concerns the retained
free kernel BEFORE any square or nilpotent filling. No survival
after such a further filling, value of the nonlinear hexagon, or
canonical fixed-lifting assertion is included.
